\documentclass[11pt]{article}
\usepackage[margin=1in]{geometry}
\usepackage{amsmath,amssymb,amsthm}
\usepackage{xurl}
\usepackage{hyperref}
\sloppy
\let\oldus\_ \renewcommand{\_}{\oldus\allowbreak}
\newtheorem{theorem}{Theorem}
\newtheorem{lemma}[theorem]{Lemma}
\theoremstyle{definition}
\newtheorem{definition}[theorem]{Definition}
\title{Conjecture 00000001046 is false:\\ $x\mapsto x+x^{q-2}$ has differential uniformity at least $4$ whenever $3\mid q-1$}
\author{Jackmeson1}
\date{2026-10-05}
\begin{document}
\maketitle

\section{The conjecture}
Conjecture 00000001046 (English text, verbatim): ``Conjecture: The differential uniformity of
$x \mapsto x + x^{q-2}$ is 2 for large $q$ (uniqueness of the optimal almost perfectly nonlinear
(APN) function for large $q$).'' The Chinese text says the same: the differential uniformity of
$x\mapsto x+x^{q-2}$ is 2 for large $q$.

\paragraph{Reading.} $q$ is a prime power and $f_q(x)=x+x^{q-2}$ is a map
$\mathbb F_q\to\mathbb F_q$. ``For large $q$'' means: there is $N$ such that
$\delta(f_q)=2$ for every $q>N$. We refute this in three settings, and claim no more:
(R1) $q$ ranges over all prime powers; (R2) $q=2^m$ (the usual APN setting); (R3) $q$ ranges over
odd prime powers. The parenthetical (uniqueness of an optimal APN function) is a separate assertion; we
neither use nor refute it. Only the explicit differential-uniformity clause is refuted, which already
makes the conjunction false. Not refuted: a reading restricted to $q=2^m$ with
$m$ odd, or to $q=3^k$ (there $3\nmid q-1$ and our argument does not apply).

\section{Definitions}
\begin{definition}
For a finite field $F$ and $g:F\to F$, the differential uniformity is
\[ \delta(g)=\max_{a\ne 0,\ b\in F}\ \#\{x\in F : g(x+a)-g(x)=b\}. \]
\end{definition}
This is the $c=1$ case of the definition quoted in the source below (Yan, arXiv:2101.10543):
``$F(x)$ is differentially $(c,\delta)$ uniform, when the maximum number of solutions
$x\in\mathrm{GF}(p^n)$ of $F(x+a)-F(cx)=b$, $a,b,c\in\mathrm{GF}(p^n)$, $c\neq1$ if $a=0$, is
equal to $\delta$.'' With $c=1$ the condition is $a\ne0$. In characteristic 2, $-$ equals $+$.

\section{Proof}
\begin{lemma}\label{lem:inv}
If $q=|F|\ge 3$ then $x^{q-2}=x^{-1}$ for $x\ne0$ and $0^{q-2}=0$. Hence $f(x)=x+x^{-1}$, with
the convention $0^{-1}=0$.
\end{lemma}
\begin{proof}
For $x\ne0$, $x^{q-1}=1$ (Lagrange in $F^\times$), so $x^{q-2}\cdot x=1$. Since $q-2\ge1$,
$0^{q-2}=0$.
\end{proof}

\begin{lemma}\label{lem:omega}
If $3\mid q-1$ there is $\omega\in F$ with $\omega\ne1$ and $\omega^3=1$; then
$\omega^2+\omega+1=0$.
\end{lemma}
\begin{proof}
$|F^\times|=q-1$ is divisible by the prime $3$, so by Cauchy's theorem $F^\times$ has an element
of order $3$. From $(\omega-1)(\omega^2+\omega+1)=\omega^3-1=0$ and $\omega\ne1$ we get
$\omega^2+\omega+1=0$.
\end{proof}

\begin{theorem}\label{thm:main}
If $3\mid q-1$ then the equation $f(x+1)-f(x)=2$ has the four distinct solutions
$x=0,-1,\omega,\omega^2$. Hence $\delta(f_q)\ge4$, so $\delta(f_q)\ne2$.
\end{theorem}
\begin{proof}
$q-1\ge3$, so $q\ge4$ and Lemma~\ref{lem:inv} applies. We have $\omega^{-1}=\omega^2$,
$(\omega^2)^{-1}=\omega$, $\omega+1=-\omega^2$ and $\omega^2+1=-\omega$.
\begin{itemize}
\item $x=0$: $f(1)-f(0)=(1+1)-0=2$.
\item $x=-1$: $f(0)-f(-1)=0-(-1-1)=2$.
\item $x=\omega$: $f(\omega+1)-f(\omega)=(-\omega^2-\omega)-(\omega+\omega^2)
  =-2(\omega^2+\omega)=2$.
\item $x=\omega^2$: $f(\omega^2+1)-f(\omega^2)=(-\omega-\omega^2)-(\omega^2+\omega)=2$.
\end{itemize}
Distinctness: $\omega\ne0$; $\omega\ne-1$ (else $\omega^2+\omega+1=1$); $\omega^2\ne-1$
(else $\omega=\omega^2+\omega+1-\omega^2-1=0$); $\omega^2\ne\omega$ (else $\omega=1$);
$0\ne-1$; $\omega^2\ne0$. Taking $a=1\ne0$, $b=2$ gives $\delta(f_q)\ge4$.
\end{proof}

\begin{theorem}
(R1) For every $N$ there is a finite field with $q>N$ and $\delta(f_q)\ge4$.
(R2) For every $k\ge1$, $\delta(f_{4^k})\ge4$; so no $N$ has $\delta(f_{2^m})=2$ for all
$m\ge N$. (R3) For every $k\ge1$, $\delta(f_{7^k})\ge4$; so no $N$ has $\delta(f_q)=2$ for all
odd $q>N$.
\end{theorem}
\begin{proof}
$4\equiv7\equiv1\pmod 3$, so $3\mid 4^k-1$ and $3\mid 7^k-1$; apply Theorem~\ref{thm:main} to
$\mathbb F_{4^k}=\mathrm{GaloisField}(2,2k)$ and $\mathbb F_{7^k}=\mathrm{GaloisField}(7,k)$,
which have $2^{2k}$ and $7^k$ elements and are arbitrarily large.
\end{proof}
Numerical check (not used in the proof): $\delta(f_{2^m})=4$ for $m=2,4,6,8$ and $=2$ for
$m=3,5,7$.

\section{Lean formalization}
File \texttt{lean/Conjecture1046/Basic.lean} (221 lines, Lean 4 + Mathlib, only import
\texttt{Mathlib}). Definitions: \texttt{f x = x + x \^{} (Fintype.card F - 2)};
\texttt{diffCount g a b} $=\#\{x\mid g(x+a)-g(x)=b\}$; \texttt{diffUniformity g} is the
\texttt{Finset.sup} of \texttt{diffCount g a b} over $a\ne0$ and all $b$. Main theorems:
\begin{itemize}
\item \texttt{four\_le\_diffUniformity}: for any finite field $F$ with
  \texttt{3 | Fintype.card F - 1}, \texttt{4 <= diffUniformity f}.
\item \texttt{not\_eventually\_two}: no $N$ such that every finite field $F$ with $N<|F|$ has
  \texttt{diffUniformity f = 2} (R1).
\item \texttt{not\_eventually\_two\_char\_two}: no $N$ such that \texttt{diffUniformity f = 2} on
  \texttt{GaloisField 2 m} for all $m\ge N$, $m\ne0$ (R2); via \texttt{char\_two\_even}.
\item \texttt{not\_eventually\_two\_odd}: the same over all finite fields of odd order larger than
  $N$ (R3); via \texttt{char\_seven}.
\end{itemize}
Supporting lemmas: \texttt{pow\_card\_sub\_two} (Lemma~\ref{lem:inv}), \texttt{exists\_cube\_root}
(Lemma~\ref{lem:omega}, from Mathlib's \texttt{exists\_prime\_orderOf\_dvd\_card}),
\texttt{four\_le\_diffCount} (Theorem~\ref{thm:main}). \texttt{\#print axioms} for each main
theorem gives \texttt{[propext, Classical.choice, Quot.sound]}; there is no \texttt{sorry} and
no \texttt{native\_decide}. Build: \texttt{lake build} with Mathlib \texttt{v4.33.1}.

\section{Relation to the earlier submission}
An earlier submission (PR \#157, orionsheep) was closed. Its Lean checked only three instances
$q=16,256,1024$, and the reviewer wrote: ``Three data points are consistent with $\delta = 2$ for
all $q > 1024$, so the asymptotic claim is untouched. The genuine argument [...] appears only in
prose. Happy to re-review a version proving the general even-$m$ case in Lean.'' This submission
proves the general case in Lean: $\delta\ge4$ for every finite field with $3\mid q-1$, which
includes $\mathbb F_{2^m}$ for every even $m\ge2$ and $\mathbb F_{7^k}$ for every $k\ge1$, and it
states the refutation of the eventual claim directly.

\section*{Source}
H.~Yan, ``On $-1$-differential uniformity of ternary APN power functions'', arXiv:2101.10543,
\url{https://arxiv.org/abs/2101.10543} (abstract; definition quoted above).
\end{document}
